% sections/H_secondary.tex: Lanes sec and demo (appendix).
% Numbers only through \R / \Rci; figures through \genfig, tables through \gentab (report/README.md).
% Sources: analysis/reports/stage2/sec_SUMMARY.md, demo_SUMMARY.md, figures/{sec,demo}/captions.json,
% tables/{sec,demo}/*, results/{sec,demo}_numbers.json; audits handoff/agent_results/syn_sec_audit_*.json.
% Figures SEC_02 and SEC_06 are placed in sections/07_secondary.tex and only referred to here.
\section{Secondary analyses in full}\label{app:sec}

This appendix reports the secondary tasks of the protocol: the instructed speed (fast, natural, slow), the
instructed grip (upright, top left, top right) and the dominant hand. It ends with a test of whether demographics
would help the hand prediction. Speed, grip and dominance are targets here. They never enter the hand model as
inputs, and no model of this appendix feeds the displayed \pright{}. Table~\ref{tab:app:sec-decisions} lists
every decision with its options, its choice and its evidence; the subsections give the detail. The main-text
summary is Section~\ref{sec:secondary}. Abbreviations used in this appendix: leave-one-person-out (LOSO),
logistic regression (LR), gradient-boosted trees (LightGBM), area under the ROC curve (AUC), confidence interval
(CI), standard deviation (SD), interquartile range (IQR), ordinary least squares (OLS), spectral arc length
(SPARC).

\begin{table}[tbp]
\centering
\small
\caption{Decisions of the secondary analyses: the options considered, the choice and where its evidence is.}
\label{tab:app:sec-decisions}
\begin{tabularx}{\linewidth}{@{}>{\raggedright\arraybackslash}p{2.5cm}>{\raggedright\arraybackslash}X%
>{\raggedright\arraybackslash}X>{\raggedright\arraybackslash}p{2.3cm}@{}}
\toprule
Decision & Options & Choice & Evidence \\
\midrule
Speed windows & the app's stored (oracle) windows; the app's detector re-run with one parameter set
  (label-blind); a fixed window from the onset & label-blind; oracle numbers kept as a label-informed,
  optimistic reference & \ref{app:sec:segmentation}, Fig.~\ref{fig:SEC_09_segmentation} \\
Common detector setting & the defaults with the longest movement cap; the slow trials' stricter end criterion
  & the defaults with the longest cap & \ref{app:sec:segmentation}, Tab.~\ref{tab:sec_segmentation} \\
Manipulation check & all trials; standard-variant trials; with or without the session position &
  standard variant, both windows, position-adjusted as sensitivity & \ref{app:sec:manipulation},
  Tab.~\ref{tab:sec_speed_manipulation} \\
Speed headline & LR on the duration; LR on 16 mirror-invariant kinematic features; LightGBM on every channel
  & the pre-specified LR on the kinematic features & \ref{app:sec:speed-models},
  Tab.~\ref{tab:sec_model_pairs} \\
Speed metrics & accuracy; balanced accuracy; macro-F1 & balanced accuracy (primary) and macro-F1, with
  per-class recall, precision and F1 & \ref{app:sec:design} \\
Speed training trials & every primary trial; standard variant only & every primary trial &
  \ref{app:sec:speed-models} \\
Speed vs the protocol & ignore; grip-only reference and grip strata & grip-only reference and grip strata &
  \ref{app:sec:confounds}, Tab.~\ref{tab:sec_design} \\
Grip model & up vector before the onset; 12 orientation features; LightGBM on every channel &
  the pre-specified 12-feature LR (identical to the up vector alone) & \ref{app:sec:grip},
  Tab.~\ref{tab:sec_grip_models} \\
Grip labels & correct the probable errors; keep as instructed & keep as instructed; explain the errors &
  \ref{app:sec:grip}, Fig.~\ref{fig:SEC_05_grip_physics} \\
Dominance inputs & every channel; mirror-invariant features only & mirror-invariant only (cannot encode
  which hand moved) & \ref{app:sec:dominance} \\
Dominance reference & AUC 0.5; within-person label-permutation null & the permutation null &
  \ref{app:sec:dominance}, Tab.~\ref{tab:sec_dominance_model} \\
Person-level handedness & a model; individual values & individual values (2 left-handed persons) &
  \ref{app:sec:dominance} \\
Demographics & add to the hand model; leave out & leave out & \ref{app:sec:demo},
  Fig.~\ref{fig:DEMO_01_paired} \\
Calibration of these models & assess; not assess & not assessed (none feeds the displayed \pright{}) &
  \ref{app:sec:design} \\
\bottomrule
\end{tabularx}
\end{table}

% ------------------------------------------------------------------------------------------------------------
\subsection{Targets, design and evaluation}\label{app:sec:design}

\paragraph{Targets.} The \R{sec/n_trials} primary trials of the \R{sec/n_persons} persons carry an instructed
speed (fast \R{sec/n_speed_fast}, natural \R{sec/n_speed_natural}, slow \R{sec/n_speed_slow}) and an instructed
grip (upright \R{sec/n_grip_upright}, top left \R{sec/n_grip_top_left}, top right \R{sec/n_grip_top_right}).
Of these, \R{sec/n_standard} are of the standard movement variant. The other variants (pause, jerky) were made at
natural speed only; \R{sec/n_fastslow_nonstandard} fast or slow trials belong to them. Slow trials are few: a
median of \R{sec/slow_per_person_median} per person, at most \R{sec/slow_per_person_max}, and
\R{sec/n_persons_no_slow} person has none.

\paragraph{Evaluation.} Every model was evaluated LOSO on the same person splits as the hand models, with the
same cap on trials per person \cite{saeb2017usecase}. Hyperparameters were tuned inside each outer fold on inner
folds of the training persons, on the class-balanced log loss, with class-balanced weights. Each model is
pre-specified: one model family whose penalty or tree parameters are tuned by nesting, so no configuration was
selected on the outer folds unless stated. Brackets are 95\,\% CIs from \R{sec/n_boot} person resamples.

\paragraph{Metrics.} The primary metric is the balanced accuracy (chance \R{sec/chance_speed} for three
classes), reported with the macro-F1 and with the recall, precision and F1 of every class. Plain accuracy would
reward answering natural, the class of \R{sec/n_speed_natural} of the \R{sec/n_trials} trials. The price of the
class-balanced weights is a low precision for the rare slow class (Appendix~\ref{app:sec:speed-models}). This is
a metric choice, not a tested remedy.

\paragraph{Paired comparisons.} Two models are compared on the trials both scored
(Tables~\ref{tab:sec_model_pairs} and~\ref{tab:sec_model_pairs_ll}). The pooled balanced-accuracy difference is
tested with an exact sign-flip test on the per-person contributions, the pooled macro-F1 difference with a
person-swap test, and the mean per-person log-loss and accuracy differences with a sign-flip test over persons.
Each metric forms its own Holm family over the \R{sec/n_model_pairs} comparisons of the two tables. A pair in which
no trial changes class and the log-loss difference is below \R{sec/ll_precision} is called identical and is not
tested.

\paragraph{No calibration claims.} None of these models feeds the displayed percentage, which is the \pright{}
of the hand model. Their log losses are reported with the other metrics; their calibration was not assessed.

\gentab{sec/sec_model_pairs}
\gentab{sec/sec_model_pairs_ll}

% ------------------------------------------------------------------------------------------------------------
\FloatBarrier
\subsection{Speed: label-blind windows}\label{app:sec:segmentation}

\paragraph{The problem.} The app resolved its online segmentation parameters per instructed condition
(Appendix~\ref{app:data:segmentation}). By default an approach ends after \R{sec/seg_default_end_sustain_ms}\,ms below
\R{sec/seg_default_end_gyro_dps}\,deg/s and \R{sec/seg_default_end_acc_ms2}\,m/s$^2$, with a movement cap of
\R{sec/seg_default_max_movement_ms}\,ms. A slow trial ends only after \R{sec/seg_slow_end_sustain_ms}\,ms below
\R{sec/seg_slow_end_gyro_dps}\,deg/s and \R{sec/seg_slow_end_acc_ms2}\,m/s$^2$, with a cap of
\R{sec/seg_slow_max_movement_ms}\,ms; pause and jerky trials also wait longer for their end. The stored (oracle)
window of a slow trial is therefore chosen by its label. Its duration, and every feature computed over it,
carries a criterion that a sensor-only system does not have.

\paragraph{Options.} (i) The oracle windows. (ii) The app's own detector re-run on every stored stream with one
parameter set for every trial: the label-blind windows. (iii) A fixed window from the label-blind onset; it was
not run in this version. For (ii), two common settings were considered: the defaults with the longest movement
cap of any condition (\R{sec/seg_common_max_movement_ms}\,ms), or the slow trials' stricter end criterion for
every trial.

\paragraph{Evidence.} With the stored parameters, the re-run detector reproduces the stored bounds of
\R{sec/reseg_n_reproduced} of \R{sec/reseg_n_trials} trials exactly. With the common default setting it keeps the
bounds of \R{sec/reseg_natural_share_std_same} of the standard natural trials and
\R{sec/reseg_fast_share_std_same} of the standard fast ones, but of \R{sec/reseg_slow_n_std_same} of the
\R{sec/reseg_slow_n_std} slow trials (Table~\ref{tab:sec_segmentation}; Figure~\ref{fig:SEC_09_segmentation}a).
The median slow approach shrinks from \R{sec/reseg_slow_dur}\,s to \R{sec/reseg_slow_dur_common}\,s, as short as
a natural one (\R{sec/reseg_natural_dur_common}\,s). It now ends after a median \R{sec/reseg_slow_path_share} of
its rotation (IQR \R{sec/reseg_slow_path_share_q25} to \R{sec/reseg_slow_path_share_q75}). The
stricter criterion is no usable common choice: applied to every trial, it ends the approach before the arm starts
back in only \R{sec/reseg_fast_strict_ok_std} of the standard fast and \R{sec/reseg_natural_strict_ok_std} of the
standard natural trials. The common default setting ends all \R{sec/reseg_n_common_ok} primary trials and
\R{sec/reseg_nat_ok} of the \R{sec/reseg_nat_n} free-choice trials. Every trial was re-standardised on the
label-blind bounds through the cache's own code path. On the oracle bounds this path reproduces the cached arrays
of all \R{sec/restd_n_checked} trials (largest difference \R{sec/restd_max_diff}).

\paragraph{Consequence.} The same pre-specified LR falls from a balanced accuracy of \Rci{sec/speed_lr_bacc} on
the oracle windows to \Rci{sec/speed_lr_common_bacc} on the label-blind windows (paired difference
\Rcib{sec/pair_speed_lr_common_vs_speed_lr_d_bacc}~points, $p < 0.001$, Holm over the balanced-accuracy family).
Its slow recall falls from \R{sec/speed_lr_rec_slow} to \R{sec/speed_lr_common_rec_slow}
(Figure~\ref{fig:SEC_09_segmentation}b). The same paired difference (label-blind minus oracle) is
\Rcib{sec/pair_speed_gbm_common_vs_speed_gbm_d_bacc}~points for LightGBM and
\Rcib{sec/pair_speed_ang_common_vs_speed_ang_d_bacc}~points for an LR on the \R{sec/speed_ang_n_features}
duration and angular features (both $p < 0.001$, Holm).

\paragraph{Decision.} The label-blind windows carry every sensor-only speed result, and the label-blind LR is the
speed headline. The oracle-window numbers are kept as a label-informed, optimistic reference. They are not an
upper bound: the label-blind numbers belong to the app's default end rule, which truncates a slow reach at its
first slow phase, and a detector built for slow movements could recover more. No segmentation parameter depends
on the hand, so the hand models are unaffected.

\begin{figure}[tbp]
\centering
\genfig{sec/SEC_09_segmentation}
\caption{The approach window depended on the instructed speed. (a) Per trial, the approach duration of the app's
stored (oracle) window against that of the app's detector re-run with one parameter set for every trial
(label-blind; dashed: equal). The median slow approach shrinks from \R{sec/reseg_slow_dur}\,s to
\R{sec/reseg_slow_dur_common}\,s; fast and natural trials barely change. (b) Recall per instructed speed of the
headline LR fitted on each kind of window (LOSO; dotted: chance). The oracle-window speed results are optimistic,
because the segmentation used the label.}
\label{fig:SEC_09_segmentation}\label{fig:app:sec-segmentation}
\end{figure}

\gentab{sec/sec_segmentation}

% ------------------------------------------------------------------------------------------------------------
\FloatBarrier
\subsection{Speed: the manipulation check}\label{app:sec:manipulation}

\paragraph{Method.} The check asks whether the persons moved differently when told to. Per person, the median
approach duration and peak angular speed were taken per instructed speed, on the \R{sec/manip_standard_n_trials}
standard-variant trials (the variants occur at natural speed only). Across persons the analysis gives the median
per-person ratio of the expected-higher to the expected-lower speed (1 = no difference) with its CI, the number of
persons in the instructed order, and a Wilcoxon signed-rank test, Holm-adjusted over the four contrasts of each
set of trials (Table~\ref{tab:sec_speed_manipulation}; Figure~\ref{fig:SEC_01_speed_manipulation}).

\paragraph{Fast vs natural.} The fast instruction raised the peak angular speed in
\R{sec/manip_standard_peak_natural_fast_n_expected} of \R{sec/manip_standard_peak_natural_fast_n} persons (ratio
\Rci{sec/manip_standard_peak_natural_fast_ratio}, $p < 0.001$, Holm). It shortened the reach in only
\R{sec/manip_standard_dur_fast_natural_n_expected} of \R{sec/manip_standard_dur_fast_natural_n} persons
(natural/fast duration ratio \Rci{sec/manip_standard_dur_fast_natural_ratio}, Holm $p$
\R{sec/manip_standard_dur_fast_natural_p_holm}). The persons followed the fast instruction at the peak; its effect
on the duration of the whole reach is weak and uncertain.

\paragraph{Order and practice.} The protocol ran the fast trials before the natural ones for
\R{sec/des_fast_before} of \R{sec/des_fast_before_n} persons, and natural reaches shortened over the session
(median within-person Spearman $\rho$ of the log duration with the position \R{sec/des_practice_rho}, negative for
\R{sec/des_practice_neg} of \R{sec/des_practice_n} persons, Wilcoxon $p$ \R{sec/des_practice_p}). The unadjusted
contrast therefore mixes instruction and practice. Two sensitivity analyses form their own Holm family with the
adjusted peak contrast. A per-person OLS of the log duration on speed and position gives natural longer than fast
for \R{sec/adj_standard_duration_log_n_expected} of \R{sec/adj_standard_duration_log_n} persons (ratio
\Rci{sec/adj_standard_duration_log_ratio}, Holm $p$ \R{sec/adj_standard_duration_log_p_holm}). Keeping only the
natural trials made before the person's last fast trial gives \R{sec/adj_standard_duration_overlap_n_expected} of
\R{sec/adj_standard_duration_overlap_n} persons (ratio \Rci{sec/adj_standard_duration_overlap_ratio}, Holm $p$
\R{sec/adj_standard_duration_overlap_p_holm}). The direction is the instructed one in each analysis; the size
stays uncertain. Adjusted for position, the peak angular speed is higher for fast in
\R{sec/adj_standard_ang_speed_peak_n_expected} of \R{sec/adj_standard_ang_speed_peak_n} persons ($p < 0.001$,
Holm).

\paragraph{Slow vs natural.} On the oracle windows, slow reaches are longer than natural ones for
\R{sec/manip_standard_dur_natural_slow_n_expected} of \R{sec/manip_standard_dur_natural_slow_n} persons (ratio
\Rci{sec/manip_standard_dur_natural_slow_ratio}, $p < 0.001$, Holm). On the label-blind windows only
\R{sec/manip_common_dur_natural_slow_n_expected} of \R{sec/manip_common_dur_natural_slow_n} are (ratio
\Rci{sec/manip_common_dur_natural_slow_ratio}, Holm $p$ \R{sec/manip_common_dur_natural_slow_p_holm}): most of the
oracle contrast is the stricter end criterion (Appendix~\ref{app:sec:segmentation}). At the peak, natural reaches
are faster than slow ones for \R{sec/manip_standard_peak_slow_natural_n_expected} of
\R{sec/manip_standard_peak_slow_natural_n} persons on the oracle windows (Holm $p$
\R{sec/manip_standard_peak_slow_natural_p_holm}) and for \R{sec/manip_common_peak_slow_natural_n_expected} of
\R{sec/manip_common_peak_slow_natural_n} on the label-blind windows (Holm $p$ \R{sec/manip_common_peak_slow_natural_p_holm}).

\paragraph{Decision.} The check uses the standard-variant trials, both kinds of window, and the position-adjusted
contrasts as a sensitivity analysis. Over every primary trial, variants included, fast reaches are shorter than
natural ones for only \R{sec/manip_all_dur_fast_natural_n_expected} of \R{sec/manip_all_dur_fast_natural_n}
persons (Holm $p$ \R{sec/manip_all_dur_fast_natural_p_holm}).

\begin{figure}[tbp]
\centering
\genfig{sec/SEC_01_speed_manipulation}
\caption{Manipulation check of the speed instruction on the oracle windows: every standard-variant trial (small
dots) and the per-person median (large dots) of (a) the approach duration and (b) the peak angular speed per
instructed speed, for each of the \R{sec/manip_standard_n_persons} persons. Panel titles: persons whose medians
follow the instruction (Wilcoxon over persons, Holm). Fast reaches were faster at the peak for
\R{sec/manip_standard_peak_natural_fast_n_expected} of \R{sec/manip_standard_peak_natural_fast_n} persons; their
shortening is uncertain and confounded with the session order. The slow durations are partly the app's stricter
end criterion.}
\label{fig:SEC_01_speed_manipulation}\label{fig:app:sec-manipulation}
\end{figure}

\gentab{sec/sec_speed_manipulation}

% ------------------------------------------------------------------------------------------------------------
\FloatBarrier
\subsection{Speed: models and paired comparisons}\label{app:sec:speed-models}

\paragraph{Options.} Four models were fitted on each kind of window (Table~\ref{tab:sec_speed_models}): a
multinomial L2 LR on the log duration alone (the baseline); the same LR on the \R{sec/n_even_features}
mirror-invariant kinematic features of Table~\ref{tab:sec_features} (the pre-specified headline); LightGBM on
summary statistics of every channel; and an LR on the up vector before the onset (the grip, a nuisance reference;
Appendix~\ref{app:sec:confounds}). Every feature is a function of one trial; scalers and models were fitted inside
the folds \cite{pedregosa2011sklearn}.

\paragraph{Headline.} On the label-blind windows the headline LR reaches a balanced accuracy of
\Rci{sec/speed_lr_common_bacc} and a macro-F1 of \Rci{sec/speed_lr_common_f1} (chance \R{sec/chance_speed}).
Recall, precision and F1 are \R{sec/speed_lr_common_rec_fast}, \R{sec/speed_lr_common_prec_fast} and
\R{sec/speed_lr_common_f1_fast} for fast; \R{sec/speed_lr_common_rec_natural},
\R{sec/speed_lr_common_prec_natural} and \R{sec/speed_lr_common_f1_natural} for natural; and
\R{sec/speed_lr_common_rec_slow}, \R{sec/speed_lr_common_prec_slow} and \R{sec/speed_lr_common_f1_slow} for slow
(Table~\ref{tab:sec_speed_classes}). The class-balanced weights make it call slow
\R{sec/speed_lr_common_npred_slow} times for \R{sec/speed_lr_common_n_slow} slow trials. The largest confusion is
between natural and slow: \R{sec/speed_lr_common_cmr_natural_slow} of the natural trials are called slow and
\R{sec/speed_lr_common_cmr_slow_natural} of the slow trials natural (Table~\ref{tab:sec_speed_confusion};
Figure~\ref{fig:SEC_02_speed_confusion}).

\paragraph{Kinematics vs duration.} The duration alone reaches \Rci{sec/speed_dur_common_bacc}. The kinematic
features add \Rcib{sec/pair_speed_lr_common_vs_speed_dur_common_d_bacc}~points of balanced accuracy ($p < 0.001$,
Holm) and \Rcib{sec/pair_speed_lr_common_vs_speed_dur_common_d_f1} of macro-F1 (Holm $p$
\R{sec/pair_speed_lr_common_vs_speed_dur_common_d_f1_p_holm}). Decision: the kinematic features, not the duration
alone.

\paragraph{LR vs LightGBM.} On the label-blind windows LightGBM reaches \Rci{sec/speed_gbm_common_bacc} (macro-F1
\Rci{sec/speed_gbm_common_f1}). Its difference from the headline is
\Rcib{sec/pair_speed_gbm_common_vs_speed_lr_common_d_bacc}~points in balanced accuracy (Holm $p$
\R{sec/pair_speed_gbm_common_vs_speed_lr_common_d_bacc_p_holm}) and
\Rcib{sec/pair_speed_gbm_common_vs_speed_lr_common_d_f1} in macro-F1 (Holm $p$
\R{sec/pair_speed_gbm_common_vs_speed_lr_common_d_f1_p_holm}); neither is significant. Its log-loss difference is
\Rcib{sec/pair_speed_gbm_common_vs_speed_lr_common_d_ll} and its accuracy difference
\Rcib{sec/pair_speed_gbm_common_vs_speed_lr_common_d_acc}~points (both $p < 0.001$, Holm;
Table~\ref{tab:sec_model_pairs_ll}). Plain accuracy weights the classes by their size, and LightGBM recalls the
largest class far more often (natural recall \R{sec/speed_gbm_common_rec_natural} against
\R{sec/speed_lr_common_rec_natural}) and slow far less often (\R{sec/speed_gbm_common_rec_slow} against
\R{sec/speed_lr_common_rec_slow}). On the oracle windows LightGBM's macro-F1 difference is
\Rcib{sec/pair_speed_gbm_vs_speed_lr_d_f1} ($p < 0.001$, Holm), while its balanced accuracy differs by
\Rcib{sec/pair_speed_gbm_vs_speed_lr_d_bacc}~points (Holm $p$ \R{sec/pair_speed_gbm_vs_speed_lr_d_bacc_p_holm}).
Decision: the LR stays the headline because it was pre-specified and is interpretable, not because it scores
best.

\paragraph{Training trials.} Trained on the standard-variant trials only, the LR reaches
\Rci{sec/speed_lr_std_bacc} on those trials. Against the all-trials model on the same
\R{sec/pair_speed_lr_std_vs_speed_lr_common_n} trials the difference is
\Rcib{sec/pair_speed_lr_std_vs_speed_lr_common_d_bacc}~points (Holm $p$
\R{sec/pair_speed_lr_std_vs_speed_lr_common_d_bacc_p_holm}). Decision: train on every primary trial.

\paragraph{Free-choice trials.} Each of the \R{sec/speed_lr_common_nat_n} free-choice trials was scored by its
person's fold model on its label-blind window. The model calls \R{sec/speed_lr_common_nat_natural} of them
natural, \R{sec/speed_lr_common_nat_fast} fast and \R{sec/speed_lr_common_nat_slow} slow. This natural-speed
recall of \R{sec/speed_lr_common_nat_share_natural} concerns free-choice trials only, all made at natural speed,
where a constant natural answer would score 100\,\%; it is not a balanced measure.

\gentab{sec/sec_speed_models}
\gentab{sec/sec_speed_classes}
\gentab{sec/sec_speed_confusion}

% ------------------------------------------------------------------------------------------------------------
\FloatBarrier
\subsection{Speed and the protocol: grip and order}\label{app:sec:confounds}

\paragraph{The confound.} The protocol confounds the instructed speed with the grip and with the position in the
session (Table~\ref{tab:sec_design}). Slow trials were upright in \R{sec/des_slow_upright} of \R{sec/des_slow_n}.
Among standard-variant trials, \R{sec/des_side_fast} of the fast, \R{sec/des_side_natural} of the natural and
\R{sec/des_side_slow} of the slow trials were held sideways. Fast trials were more often sideways than natural
ones for \R{sec/des_fast_more_side} of \R{sec/des_fast_more_side_n} persons. The median relative position in the
session (0 = first, 1 = last) was \R{sec/des_pos_slow} for slow, \R{sec/des_pos_fast} for fast and
\R{sec/des_pos_natural} for natural trials; slow came before natural for \R{sec/des_slow_before} of
\R{sec/des_slow_before_n} persons. A model can therefore predict the speed from the grip or from practice instead
of from the movement.

\paragraph{Grip-only reference.} An LR on the up vector before the onset alone predicts the instructed speed with a
balanced accuracy of \Rci{sec/speed_grip_bacc} (slow recall \R{sec/speed_grip_rec_slow}). The headline does not
differ from it significantly in balanced accuracy (\Rcib{sec/pair_speed_lr_common_vs_speed_grip_d_bacc}~points,
Holm $p$ \R{sec/pair_speed_lr_common_vs_speed_grip_d_bacc_p_holm}) but beats it in macro-F1
(\Rcib{sec/pair_speed_lr_common_vs_speed_grip_d_f1}, $p < 0.001$, Holm). The mirror-invariant features themselves
carry some grip information: they separate sideways from upright trials with a balanced accuracy of
\Rci{sec/side_even_bacc} (chance \R{sec/chance_side}).

\paragraph{Within grip strata.} The models were re-scored within the upright and within the sideways trials after
prediction (Tables~\ref{tab:sec_speed_strata} and~\ref{tab:sec_speed_strata_pairs}). A speed enters a stratum only
with trials of at least \R{sec/str_min_persons} persons, so slow is left out of the sideways stratum. Within the
upright trials the headline reaches \R{sec/str_speed_lr_common_upright_bacc} over three classes and LightGBM
\R{sec/str_speed_gbm_common_upright_bacc} (difference
\Rcib{sec/str_pair_speed_gbm_common_vs_speed_lr_common_upright_d_bacc}~points, Holm $p$
\R{sec/str_pair_speed_gbm_common_vs_speed_lr_common_upright_d_bacc_p_holm}). Within the sideways trials (fast and
natural) they reach \R{sec/str_speed_lr_common_sideways_bacc} and \R{sec/str_speed_gbm_common_sideways_bacc}
(\Rcib{sec/str_pair_speed_gbm_common_vs_speed_lr_common_sideways_d_bacc}~points, Holm $p$
\R{sec/str_pair_speed_gbm_common_vs_speed_lr_common_sideways_d_bacc_p_holm}). LightGBM, which sees the
device-frame channels, predicts fast for \R{sec/str_speed_gbm_common_upright_pred_fast} of the upright and
\R{sec/str_speed_gbm_common_sideways_pred_fast} of the sideways trials, close to the instructed shares
(\R{sec/str_speed_gbm_common_upright_true_fast} and \R{sec/str_speed_gbm_common_sideways_true_fast}). The headline
predicts fast for \R{sec/str_speed_lr_common_upright_pred_fast} and \R{sec/str_speed_lr_common_sideways_pred_fast}.

\paragraph{Decision.} The grip-only model is reported as a nuisance reference, and the model comparisons are
repeated within grip strata. The speed accuracies may still partly reflect the protocol's grip and order
structure.

\gentab{sec/sec_design}
\gentab{sec/sec_speed_strata}
\gentab{sec/sec_speed_strata_pairs}

% ------------------------------------------------------------------------------------------------------------
\FloatBarrier
\subsection{Speed: which features carry it}\label{app:sec:speed-features}

\paragraph{Model-free effects.} Per person, the medians of each kinematic feature were compared between two
instructed speeds on the standard-variant trials, in units of the feature's SD over all trials
(Table~\ref{tab:sec_speed_univariate}; Figure~\ref{fig:SEC_03_speed_features}a). A Wilcoxon test over persons is
Holm-adjusted over every feature and both contrasts (\R{sec/n_even_features} features, two contrasts). On the
label-blind windows, fast differs from natural most in \R{sec/uni_common_fast_top1_feature}
(\R{sec/uni_common_fast_top1_d}\,SD, positive for \R{sec/uni_common_fast_top1_npos} of
\R{sec/uni_common_fast_top1_n} persons, Holm $p$ \R{sec/uni_common_fast_top1_p_holm}),
\R{sec/uni_common_fast_top2_feature} (\R{sec/uni_common_fast_top2_d}\,SD, \R{sec/uni_common_fast_top2_npos} of
\R{sec/uni_common_fast_top2_n}) and \R{sec/uni_common_fast_top3_feature} (\R{sec/uni_common_fast_top3_d}\,SD,
\R{sec/uni_common_fast_top3_npos} of \R{sec/uni_common_fast_top3_n}). In all, \R{sec/uni_common_fast_n_sig}
features differ significantly between fast and natural, and \R{sec/uni_common_slow_n_sig} between slow and natural. The
largest slow effect is in \R{sec/uni_common_slow_top1_feature} (\R{sec/uni_common_slow_top1_d}\,SD, Holm $p$
\R{sec/uni_common_slow_top1_p_holm}). On the oracle windows the counts are \R{sec/uni_oracle_fast_n_sig} and
\R{sec/uni_oracle_slow_n_sig}, led by the slow trials' \R{sec/uni_oracle_slow_top1_feature}
(\R{sec/uni_oracle_slow_top1_d}\,SD): the oracle slow contrasts are largely the label-dependent end criterion.

\paragraph{Model coefficients.} The headline LR's largest coefficients (Figure~\ref{fig:SEC_03_speed_features}b)
are \R{sec/coef_fast_top1_feature} (\R{sec/coef_fast_top1} log-odds per SD) and \R{sec/coef_fast_top2_feature}
(\R{sec/coef_fast_top2}) for fast vs natural, and \R{sec/coef_slow_top1_feature} (\R{sec/coef_slow_top1}) and
\R{sec/coef_slow_top2_feature} (\R{sec/coef_slow_top2}) for slow vs natural. These coefficients are conditional on
the other, correlated features and vary across folds, so the model-free effects are the better guide to what
carries the speed. Table~\ref{tab:sec_features} defines every feature and its mirror parity; smoothness follows
the SPARC and jerk measures of the literature \cite{melendezcalderon2021smoothness,flash1985coordination}.

\begin{figure}[tbp]
\centering
\genfig{sec/SEC_03_speed_features}
\caption{Which kinematic features carry the instructed speed, on label-blind windows. (a) Per person, the
difference of the medians between two instructed speeds in units of the feature's SD (median and interquartile
range over persons; standard-variant trials). (b) The log-odds contrasts of the headline LR per SD of each
feature (mean and range over the outer folds; conditional on the other features). Fast differs from natural most
in \R{sec/uni_common_fast_top1_feature} and \R{sec/uni_common_fast_top2_feature}; no slow effect survives Holm.}
\label{fig:SEC_03_speed_features}\label{fig:app:sec-features}
\end{figure}

% The generated table sets \small and is slightly taller than a page at that size; typeset it one size smaller.
{\let\small\footnotesize\renewcommand{\arraystretch}{0.95}\gentab{sec/sec_speed_univariate}}
\gentab{sec/sec_features}

% ------------------------------------------------------------------------------------------------------------
\FloatBarrier
\subsection{Speed: who follows the instruction}\label{app:sec:followers}

Figure~\ref{fig:SEC_08_speed_followers} is a description, not evidence. Per person it plots the natural/fast
duration ratio on the label-blind windows against the headline model's two-class fast-vs-natural balanced
accuracy on the same standard trials (the class decided between fast and natural only; chance 0.5). The two are
correlated (Spearman $\rho$ \R{sec/fol_rho}, $p$ \R{sec/fol_p}, \R{sec/fol_n} persons, one exploratory test). The
mean two-class balanced accuracy is \R{sec/fol_bacc_followed} for the \R{sec/fol_n_followed} persons whose fast
reaches were shorter and \R{sec/fol_bacc_not} for the other \R{sec/fol_n_not}. The analysis is partly circular,
because the duration it stratifies by is an input of the model it scores, and it was added after seeing the
results.

\begin{figure}[tbp]
\centering
\genfig{sec/SEC_08_speed_followers}
\caption{Description, not evidence: per person, the median natural/fast duration ratio (standard-variant trials,
label-blind windows; right of the dashed line: fast reaches shorter) against the headline model's two-class
fast-vs-natural balanced accuracy on the same trials (dotted: chance). Open circles: the persons whose fast
reaches were not shorter. Partly circular, because the duration is an input of the model.}
\label{fig:SEC_08_speed_followers}\label{fig:app:sec-followers}
\end{figure}

% ------------------------------------------------------------------------------------------------------------
\FloatBarrier
\subsection{Grip: the up vector and the errors}\label{app:sec:grip}

\paragraph{Physics.} The grip is the orientation of the phone in the hand, so it shows in where gravity points in
the device frame during the still hold before the onset (Figure~\ref{fig:SEC_05_grip_physics}). The median unit up
vector $\ghat$ (device $x$, $y$, $z$) is (\R{sec/upmed_upright_x}, \R{sec/upmed_upright_y},
\R{sec/upmed_upright_z}) for upright, with the top edge up; (\R{sec/upmed_top_left_x}, \R{sec/upmed_top_left_y},
\R{sec/upmed_top_left_z}) for top left, with the right edge up; and (\R{sec/upmed_top_right_x},
\R{sec/upmed_top_right_y}, \R{sec/upmed_top_right_z}) for top right, with the left edge up. The device-frame
channels carry the grip directly. The gravity-frame channels do not encode the orientation, but the grip changes
how the arm moves (Appendix~\ref{app:sec:confounds}).

\paragraph{Models.} Three nested models were fitted on the hold before the onset (Tables~\ref{tab:sec_grip_models}
and~\ref{tab:sec_grip_classes}; Figure~\ref{fig:SEC_04_grip_confusion}): an LR on the up vector alone, the
pre-specified headline LR on the \R{sec/n_grip_features} device-frame orientation features, and LightGBM on the
summary statistics of every channel. The headline reaches a balanced accuracy of \Rci{sec/grip_lr_bacc} against
the instruction (chance \R{sec/chance_grip}), with recalls of \R{sec/grip_lr_rec_upright} for upright,
\R{sec/grip_lr_rec_top_left} for top left and \R{sec/grip_lr_rec_top_right} for top right. The up-vector model makes
the same predictions: \R{sec/pair_grip_lr_vs_grip_up_n_changed} trials change class and the log-loss difference is
\R{sec/pair_grip_lr_vs_grip_up_d_ll}, so the two models are identical and are not tested. LightGBM reaches
\Rci{sec/grip_gbm_bacc}; its difference from the headline is \Rcib{sec/pair_grip_gbm_vs_grip_lr_d_bacc}~points
(Holm $p$ \R{sec/pair_grip_gbm_vs_grip_lr_d_bacc_p_holm}; Table~\ref{tab:sec_model_pairs}). Decision: the headline
stays the pre-specified 12-feature LR; the up vector before the onset carries all of its information.

\paragraph{Errors.} The headline misclassifies \R{sec/grip_lr_n_errors} of \R{sec/n_trials} trials, all from
\R{sec/grip_lr_n_persons_with_errors} persons (Tables~\ref{tab:sec_grip_confusion} and~\ref{tab:sec_grip_errors}):
\R{sec/grip_err1_person} (\R{sec/grip_err1_pattern}), \R{sec/grip_err2_person} (\R{sec/grip_err2_pattern}) and
\R{sec/grip_err3_person} (\R{sec/grip_err3_pattern}). Of the errors, \R{sec/grip_lr_n_sideways_swaps} are swaps
between top left and top right. In \R{sec/grip_lr_n_up_matches_predicted} of the \R{sec/grip_lr_n_errors} errors
the up vector lies closer to the median of the predicted grip than to that of the instructed one (ringed in
Figure~\ref{fig:SEC_05_grip_physics}). A sign error is unlikely: the gravity sign was inferred from the instructed
orientation for \R{sec/grip_sign_inferred} trials, and \R{sec/grip_lr_err_upright_correct} of the
\R{sec/grip_lr_err_upright_n} upright trials of these persons are recognised. These trials were most likely held in
the predicted grip.

\paragraph{Consistency checks.} On the \R{sec/grip_lr_ref_n_confirmed} trials that neither the app nor the gravity
rule flagged, the headline's balanced accuracy is \R{sec/grip_lr_ref_bacc_confirmed}. Its prediction equals the
app's own start grip for \R{sec/grip_lr_ref_n_agree} of \R{sec/grip_lr_ref_n_known} trials
(Table~\ref{tab:sec_grip_reference}). Both references read the same up vector as the model, so they agree by
construction: they are consistency checks, not an independent confirmation. No video or instruction log records
the grip actually held.

\paragraph{Decision.} The grip labels are kept as instructed; the errors are explained by physics, not corrected.
Any result stratified by instructed grip, in every lane, mis-assigns these \R{sec/grip_lr_n_errors} trials.

\begin{figure}[tbp]
\centering
\genfig{sec/SEC_04_grip_confusion}
\caption{Confusion matrices of the instructed grip (LOSO; counts and row shares; rows: instructed grip) for the
up vector before the onset alone, the headline orientation LR and LightGBM, with their balanced accuracies. The
first two are identical. The headline's \R{sec/grip_lr_n_errors} errors include \R{sec/grip_lr_n_sideways_swaps}
swaps between top left and top right.}
\label{fig:SEC_04_grip_confusion}\label{fig:app:sec-grip-confusion}
\end{figure}

\begin{figure}[tbp]
\centering
\genfig{sec/SEC_05_grip_physics}
\caption{Why the grip is recognised and where it fails: the unit up vector in the device frame during the still
hold before the onset, per trial, coloured by the instructed grip; (a) $x$ against $y$, (b) $x$ per instructed grip
and hand. Upright puts the top edge ($+y$) up, top left the right edge ($+x$), top right the left edge ($-x$).
Ringed: the trials the headline misclassifies; each of the \R{sec/grip_lr_n_up_matches_predicted} lies in the
cluster of the predicted grip.}
\label{fig:SEC_05_grip_physics}\label{fig:app:sec-grip-physics}
\end{figure}

\gentab{sec/sec_grip_models}
\gentab{sec/sec_grip_classes}
\gentab{sec/sec_grip_confusion}
\gentab{sec/sec_grip_errors}
\gentab{sec/sec_grip_reference}

% ------------------------------------------------------------------------------------------------------------
\FloatBarrier
\subsection{Handedness and the dominant hand}\label{app:sec:dominance}

\paragraph{Question.} Can a LOSO model tell a dominant-hand trial from mirror-invariant kinematics alone? These
inputs are identical for a trial and its mirror image (Section~\ref{sec:physics}, Eq.~\eqref{eq:mirror}): every
feature was bit-identical for \R{sec/mirror_n_identical} of \R{sec/mirror_n_trials} trials and their mirror images
(largest difference \R{sec/mirror_max_abs_diff}). They cannot encode which hand moved, only how the dominant hand
moves \cite{sainburg2002dominance,przybyla2012dominance}. The dominant hand is the right hand for the
\R{sec/dom_n_right} right-handed persons and the left hand for the \R{sec/dom_n_left} left-handed ones, so
dominance is confounded with the right hand for the majority. Every person has at least
\R{sec/bal_min_dominant} dominant and \R{sec/bal_min_non_dominant} non-dominant trials. The first block was
left-handed in \R{sec/bal_first_left} sessions and right-handed in \R{sec/bal_first_right} (exact binomial $p$
\R{sec/bal_first_p}).

\paragraph{Reference.} The pre-specified reference is a label-permutation null, not an AUC of 0.5. The labels
were shuffled within each person \R{sec/dom_n_perm} times and the whole nested pipeline was re-run each time,
which keeps each person's class balance.

\paragraph{Result.} The pre-specified LR reaches a pooled AUC of \Rci{sec/dom_auc}, against a null mean of
\R{sec/dom_null_pooled_mean} (95th percentile \R{sec/dom_null_pooled_q95}; one-sided permutation $p$
\R{sec/dom_p_pooled}; Table~\ref{tab:sec_dominance_model}; Figure~\ref{fig:SEC_06_dominance_auc}). The mean
within-person AUC is \Rci{sec/dom_within_auc} (null mean \R{sec/dom_null_within_mean}, $p$ \R{sec/dom_p_within});
it is above 0.5 for \R{sec/dom_n_auc_above_half} of \R{sec/dom_n_persons} persons. The accuracy is
\Rci{sec/dom_acc} and the log loss \R{sec/dom_ll}, against \R{sec/dom_ll_const} for a constant 0.5. The nested
tuning chose the smallest penalty parameter of the grid (\R{sec/dom_c_min}, the strongest regularisation) in
\R{sec/dom_n_folds_grid_min} of \R{sec/dom_n_folds} outer folds, and the predicted probability spans only
\R{sec/dom_p_min} to \R{sec/dom_p_max}. No dominant-hand signal is detected; the CI still allows a pooled AUC up to
\R{sec/dom_auc_hi}.

\paragraph{Exploratory second model.} LightGBM (\R{sec/domgbm_n_configs} model classes tried in all, no permutation
baseline) reaches a pooled AUC of \Rci{sec/domgbm_auc}, above the LR's null mean but close to its 95th percentile.
Against the LR on the same trials, the per-person accuracy differs by \Rcib{sec/pair_dom_d_acc}~points (sign-flip
$p$ \R{sec/pair_dom_d_acc_p}) and the log loss by \Rci{sec/pair_dom_d_ll} ($p$ \R{sec/pair_dom_d_ll_p}; one
exploratory comparison, unadjusted). A weak signal that only a non-linear model finds cannot be excluded.

\paragraph{Per person.} Table~\ref{tab:sec_dominance_persons} gives each person's within-person AUC and accuracy.
The left-handed persons are shown individually, without a CI: \R{sec/dom_lh1_person} \R{sec/dom_lh1_auc} and
\R{sec/dom_lh2_person} \R{sec/dom_lh2_auc}, against a right-handed median of \R{sec/dom_rh_auc_median} (range
\R{sec/dom_rh_auc_min} to \R{sec/dom_rh_auc_max}).

\paragraph{Within-person comparison.} Per person, the medians of six kinematic features were compared between the
dominant and the non-dominant hand (Table~\ref{tab:sec_dominance_paired}; Figure~\ref{fig:SEC_07_dominance_paired};
Wilcoxon over persons, Holm over the \R{sec/pd_n_features} features). No feature differs detectably:
\R{sec/pd_n_sig} have a Holm $p$ below 0.05, and none does when matched on speed, grip and variant. The CIs bound
the differences still compatible with the data: the duration changes by \R{sec/pd_duration_log_med}~[\Rb{sec/pd_duration_log_lo},
\Rb{sec/pd_duration_log_hi}] (dominant higher for \R{sec/pd_duration_log_npos} of \R{sec/pd_duration_log_n}
persons), the peak angular speed by \R{sec/pd_ang_speed_peak_med}~[\Rb{sec/pd_ang_speed_peak_lo},
\Rb{sec/pd_ang_speed_peak_hi}], the angular path length by \R{sec/pd_ang_path_med}~[\Rb{sec/pd_ang_path_lo},
\Rb{sec/pd_ang_path_hi}] (\R{sec/pd_ang_path_npos} of \R{sec/pd_ang_path_n} higher; Holm $p$
\R{sec/pd_ang_path_p_holm}) and the linear path length by \R{sec/pd_lin_path_med}~[\Rb{sec/pd_lin_path_lo},
\Rb{sec/pd_lin_path_hi}]. SPARC differs by \R{sec/pd_ang_sparc_med}~[\Rb{sec/pd_ang_sparc_lo},
\Rb{sec/pd_ang_sparc_hi}] and the log dimensionless jerk by \R{sec/pd_ang_ldlj_med}~[\Rb{sec/pd_ang_ldlj_lo},
\Rb{sec/pd_ang_ldlj_hi}]. The left-handed persons are shown individually (diamonds); their duration differences
are \R{sec/pd_lh1_duration_log} (\R{sec/dom_lh1_person}) and \R{sec/pd_lh2_duration_log} (\R{sec/dom_lh2_person}).

\paragraph{Concordance.} Table~\ref{tab:sec_concordance} counts the persons by self-reported handedness, usual phone
hand and the hand of the first free-choice trial; it is descriptive only. Of the \R{sec/conc_n_phone_clear}
persons who name one usual phone hand, \R{sec/conc_n_phone_right} name the right hand, and \R{sec/conc_nat_right} of the \R{sec/conc_n_nat} first
free-choice trials used the right hand, which a constant would also give. The only contrast comes from the
\R{sec/conc_n_left} left-handed persons: \R{sec/conc_left_nat_right} of them chose the right hand for the free trial
and \R{sec/conc_left_phone_right} usually hold the phone in the right hand. Whether handedness or the usual phone
hand goes with the free choice cannot be told apart with so few left-handed persons and no left-hand first free
choice \cite{papadatoupastou2020handedness}. Of the \R{sec/conc_rep_n} free-choice trials of repeat visits,
\R{sec/conc_rep_left} used the left hand.

\paragraph{Decision.} No person-level handedness model is trained or evaluated: the \R{sec/dom_n_left} left-handed
persons are fewer than the \R{sec/min_persons_ci} persons a stratum needs for a CI, so their values are shown
individually. Self-reported handedness is the only dominance measure; no handedness questionnaire was given.

\begin{figure}[tbp]
\centering
\genfig{sec/SEC_07_dominance_paired}
\caption{Dominant vs non-dominant movements within each person: per person, the difference of the medians
(dominant minus non-dominant; a percent change for the duration, the peak angular speed and the path lengths);
the bar is the median over persons. Wilcoxon over persons, Holm over the \R{sec/pd_n_features} features. Diamonds:
the \R{sec/dom_n_left} left-handed persons, shown individually. No feature differs detectably; the CIs still allow
moderate differences (text).}
\label{fig:SEC_07_dominance_paired}\label{fig:app:sec-dominance-paired}
\end{figure}

\gentab{sec/sec_dominance_model}
\gentab{sec/sec_dominance_persons}
\gentab{sec/sec_dominance_paired}
\gentab{sec/sec_concordance}

% ------------------------------------------------------------------------------------------------------------
\FloatBarrier
\subsection{Demographics add-on test}\label{app:sec:demo}

This test is a quick first pass and has not been audited. Production does not know a user's age, sex, height,
handedness or usual phone hand; the test asks whether these would help if it did. Nothing from it enters the
production model.

\paragraph{Design.} Refitting the final model with covariate inputs was too slow for this pass. Instead, a
second-stage L2 LR was fitted on the final model's logit (Appendix~\ref{app:final}), with and without the person
covariates: age, sex, height (with a missing indicator; \R{demo/profile_n_height_missing} heights are missing and
were median-imputed inside the fold), self-reported handedness and usual phone hand. In each outer fold the stage
was fitted on the final model's inner cross-fitted logits of the training persons and applied to the held-out
person's logit from the saved LOSO predictions; it never uses other persons' outer-fold predictions. Its penalty
was tuned on grouped inner folds. The variants are the sensor-only stage (the comparator), the stage with the
covariates as main effects, the stage with covariates and covariate $\times$ logit interactions, the covariates
alone (the floor), and two fixed rules for the base-rate note: answer the usual phone hand, or always answer right.
Paired differences against the sensor-only stage are mean per-person differences with paired CIs and sign-flip
$p$-values, Holm-adjusted over the two add-on variants within each metric.

\paragraph{Cohort.} The \R{demo/profile_n_persons} persons include \R{demo/profile_n_left_handed} left-handed
persons, \R{demo/profile_n_phone_both} who use both hands for the phone and \R{demo/profile_n_phone_left} who
usually use the left; \R{demo/profile_n_female} are female and the ages span \R{demo/profile_age_min} to
\R{demo/profile_age_max} years. Interactions such as `left-handed persons move differently' are therefore barely
identified.

\paragraph{Results.} Table~\ref{tab:demo_models} and Figure~\ref{fig:DEMO_01_paired} give the results. The
sensor-only stage is neutral against the final model as shipped (accuracy \R{demo/stack_vs_shipped_dacc}~points,
log loss \R{demo/stack_vs_shipped_dll}; unadjusted sign-flip $p$ \R{demo/stack_vs_shipped_p_acc} and
\R{demo/stack_vs_shipped_p_ll}) and reaches an accuracy of \Rci{demo/sensor_acc}. Adding the covariates gives
\Rci{demo/main_acc}: a paired difference of \Rci{demo/main_dacc}~points (Holm $p$ \R{demo/main_holm_acc}) and a
log-loss difference of \Rci{demo/main_dll} (Holm $p$ \R{demo/main_holm_ll}). The covariates change
\R{demo/main_flips} primary answers; \R{demo/main_better} person gains and \R{demo/main_worse} lose. With the interactions the accuracy is
\Rci{demo/inter_acc} (difference \Rci{demo/inter_dacc}~points, Holm $p$ \R{demo/inter_holm_acc}), and the log loss
rises by \Rci{demo/inter_dll} (Holm $p$ \R{demo/inter_holm_ll}, significant). The interactions change
\R{demo/inter_flips} answers; \R{demo/inter_better} persons gain and \R{demo/inter_worse} lose. The covariates-only floor sits at chance:
accuracy \Rci{demo/covonly_acc}, balanced accuracy \R{demo/covonly_bacc}, AUC \R{demo/covonly_auc}. This is
expected, since every person did both hands: a person-level covariate can help only through interactions, and the
interaction variant does not help.

\paragraph{Exploratory base.} The absolute levels inherit the optimism of the final model as shipped, which is
exploratory: the best of \R{m3/deep_ship_n_configs} configurations in lane m3, then chosen among
\R{final/n_candidates} candidates. Its nested, selection-aware accuracy is \Rci{final/nested_accuracy}
(Appendix~\ref{app:final}). The paired differences share the base, so its optimism largely cancels in them. The
stage's penalty is chosen on the same inner logits it is fitted on, a second-order effect that is not removed.

\paragraph{Natural use and the base rate.} On natural use the measure is the right-hand recall on the free-choice
trials (all right-hand so far; a constant right answer would score 100\,\%). The shipped model and the sensor-only
stage reach \R{demo/shipped_nat_recall} and \R{demo/sensor_nat_recall}, the two add-ons \R{demo/main_nat_recall}
and \R{demo/inter_nat_recall}. The fixed rules show why this measure cannot judge an add-on
(Figure~\ref{fig:DEMO_02_base_rate}). Always answering right scores \R{demo/ruleright_nat_recall} and the
phone-hand rule \R{demo/rulephone_nat_recall}, yet both sit at chance on the primary set (balanced accuracy
\R{demo/ruleright_bacc} and \R{demo/rulephone_bacc}). The phone-hand rule abstains (\pright{} exactly 0.5) for the
\R{demo/profile_n_phone_both} persons who use both hands. On natural use it answers \R{demo/rulephone_nat_right} of
the \R{demo/nat_n} trials, all correctly; its score gives each of the \R{demo/rulephone_nat_ties} abstentions half
credit, and it has no forced choice for them. The floor's \R{demo/covonly_nat_recall} is not a signal: its
\pright{} never moves more than \R{demo/covonly_max_dev} from 0.5 and sits just below it, because
\R{demo/primary_right_share} of the primary trials are right-hand trials.

\paragraph{Decision.} Demographics are not added, even as a hypothetical: no add-on variant improves accuracy or
log loss on the same folds. With \R{demo/profile_n_left_handed} left-handed persons and
\R{demo/profile_n_phone_left} who usually hold the phone in the left hand, this is a null result on a small, mostly
right-handed sample, not evidence of no effect. Not done in this pass: a refit of the network with covariate
inputs, a replication on a cheaper base model, a calibration analysis of the add-ons, and an audit.

\begin{figure}[tbp]
\centering
\genfig{demo/DEMO_01_paired}
\caption{Demographics add-on test (quick first pass, not audited): the covariates added to a second stage on the
final model's logit, against the same stage without them, on the same LOSO folds. Mean per-person differences with
paired 95\,\% CIs and Holm-adjusted sign-flip $p$ over the two add-on variants within each metric (filled red:
below 0.05). Grey: the sensor-only stage against the shipped model (unadjusted). Negative accuracy and positive
log-loss differences mean the covariates hurt.}
\label{fig:DEMO_01_paired}\label{fig:app:sec-demo-paired}
\end{figure}

\begin{figure}[tbp]
\centering
\genfig{demo/DEMO_02_base_rate}
\caption{The base-rate effect: balanced accuracy on the primary set (both hands for every person) and right-hand
recall on the \R{demo/nat_n} free-choice trials (all right-hand so far; a constant right answer would score
100\,\%). The fixed rules look excellent on natural use and sit at chance on the primary set, so natural use cannot
judge an add-on; the add-on test is judged on the primary set.}
\label{fig:DEMO_02_base_rate}\label{fig:app:sec-demo-base-rate}
\end{figure}

\gentab{demo/demo_models}

% ------------------------------------------------------------------------------------------------------------
\FloatBarrier
\subsection{Caveats and open questions}\label{app:sec:caveats}

\begin{itemize}[nosep]
\item None of these models runs in production. A production speed or grip estimate would take its bounds from the
  cue-free detector (Appendix~\ref{app:prod}), whose end detection meets the same difficulty with slow reaches; its
  sensitivity to end-detection errors would have to be measured.
\item The label-blind speed estimate belongs to one detector setting, the app's defaults. A fixed window from the
  onset or a detector tuned for slow movements would give other numbers; neither was run.
\item The linear path length integrates the acceleration with zero velocity at both bounds. The drift that the
  linear correction cannot remove grows with the duration.
\item Slow trials are few (a median of \R{sec/slow_per_person_median} per person), so per-person slow statistics
  rest on very few movements.
\item Added after seeing the results (exploratory; $p$-values not corrected for that): the per-person description
  of Figure~\ref{fig:SEC_08_speed_followers}, the LightGBM dominance model, the grip consistency checks, the
  segmentation check, the grip-only reference, the grip strata and the position-adjusted contrasts.
\item The null dominance result does not test whether the hand model's evidence lies only in the chiral channels;
  that remains a hypothesis.
\item Open: a person-relative speed measure, normalised by the person's own natural movements; and a
  dominance-related movement style that transfers to left-handed persons, which needs more left-handed persons.
\end{itemize}
